% ===================================================================
%  TABELLA N-RO 5 — Le casa e su construction.
%  Traduction 2026 d'apres texto/io, controlee sur texto/fr et
%  texto/es. Consigne : voir texto/ia/10-tabella-01.tex.
%
%  Le tableau 5 est un DIALOGUE, et les deux volumes composent le nom
%  du parleur en petites capitales : on garde \textsc{}.
%
%  LES DEUX NOMS PROPRES SE TRADUISENT, ET L'IDO L'A DEJA FAIT :
%  Leblanc est « Blanko », Roux est « Rufo ». L'interlingua forme les
%  memes mots sur les memes adjectifs — blanc, rufe — et ecrit donc
%  senior Blanco et senior Rufo.
% ===================================================================

%%K t05-apar-1 apar
\VUsaut{6.0mm}
\VUtitre{15.0pt}{60}{TABELLA N\textsuperscript{o} 5}
\VUsaut{1.4mm}
\VUtitre{11.4pt}{0}{Le casa e su construction.}
\VUsaut{2.2mm}

%%K t05-tit-1 sub
\VUpk{10.2pt}{40}{Bon incontro.}
\VUsaut{3.0mm}

%%K t05-01-1 p
\textsc{Johannes}. --- Oncle mie, io desira multo saper como on construe un
\VUgras{casa}.

%%K t05-01-2 p
\textsc{Le oncle} (80). --- Isto es multo facile, mi amico, veni con me, io te
monstrara illo que senior Bernardo (79) face construer e que io habitara.

%%K t05-01-3 p
\textsc{Johannes}. --- Qui dunque designava le \VUgras{plano}
\textsuperscript{(76)} de iste casa?

%%K t05-01-4 p
\textsc{Le oncle}. --- \VUgras{Le architecto} \textsuperscript{(75)}; ille
presentava un multo clar designo que esseva approbate, e \VUgras{le
interprenditor} \textsuperscript{(77)} comenciava le labores.

%%K t05-01-5 p
\textsc{Johannes}. --- Qui dunque es le interprenditor?

%%K t05-01-6 p
\textsc{Le oncle}. --- Ille es senior Blanco, le patre de tu amico Jacobo. Ille
dirige le obreros con senior Rufo, le architecto.

%%K t05-01-7 p
Ecce! justo a tempore veni senior Blanco con su \VUgras{regula}
\textsuperscript{(78)}, e senior Rufo con su plano.

%%K t05-01-8 p
Bon die, seniores. Como sta vos?

%%K t05-01-9 p
\textsc{Senior Rufo}. --- Multo ben, gratias. Bon die, Johannes, quanto cresceva
iste infante!

%%K t05-01-10 p
\textsc{Le oncle}. --- Si, vermente, ille es un parve homine. Ille es multo
seriose, on debe informar le de toto lo que ille vide. Hodie, ille volerea
saper como on construe un casa.

%%K t05-tit-2 sub
\VUpk{10.2pt}{40}{Le obreros.}
\VUsaut{3.0mm}

%%K t05-01-11 p
\textsc{Senior Blanco}. --- Ben, veni con me, parve amico, io te lo explicara
tosto, proque io ama multo le pueros qui vole instruer se.

%%K t05-01-12 p
Tu vide ille obrero qui tene un \VUgras{pala} \textsuperscript{(82)} in le mano
: ille es un \VUgras{terrassero} \textsuperscript{(81)}. Ille excava un
\VUgras{trenchea} \textsuperscript{(46)} pro poner le tubos de aqua. Ille anque
excavava le solo in le loco ubi es le casa, pro que le muratero erige le quatro
\VUgras{muros}. Le parte de iste muros que es in le solo es nominate
\VUgras{fundamentos} \textsuperscript{(2)}. Le spatio inserrate inter le
fundamentos forma le \VUgras{cellario} \textsuperscript{(3)}. Iste cellario ha
un parve fenestra nominate \VUgras{fenestra de cellario}
\textsuperscript{(5)}, un \VUgras{porta} \textsuperscript{(4)} e un scala.

%%K t05-01-13 p
Le \VUgras{muratero} \textsuperscript{(108)} continuava construer le muros
lassante aperturas pro le \VUgras{portas} \textsuperscript{(8)} e le
\VUgras{fenestras} \textsuperscript{(17)}.

%%K t05-01-14 p
\textsc{Johannes}. --- Ma, iste muratero io non le vide!

%%K t05-01-15 p
S\textsuperscript{o} \textsc{Blanco}. --- Ora ille finiva laborar in le casa.
Ille es presso le \VUgras{stabulo} \textsuperscript{(54)}, sur su
\VUgras{scafalage} \textsuperscript{(110)}, ille construe le muro del
\VUgras{lavanderia} \textsuperscript{(56)}.

%%K t05-01-16 p
Ille ha un trulla, un gaveta e un \VUgras{fil a plumbo} \textsuperscript{(109)}.
Ma tu non memorara iste nomines.

%%K t05-01-17 p
\textsc{Johannes}. --- Certo si, proque io tosto demandara a mi oncle un parve
trulla e un gaveta, e io mesme fabricara un fil a plumbo con un cordetta.

%%K t05-tit-3 sub
\VUsaut{3.0mm}
\VUpk{10.2pt}{40}{Le material.}
\VUsaut{3.0mm}

%%K t05-01-18 p
\textsc{Le oncle}. --- Ha! multo ben. Ma dice me, Johannes, que ha on besonio
pro construer un casa?

%%K t05-01-19 p
\textsc{Johannes}. --- On ha besonio de \VUgras{petras de talia}
\textsuperscript{(59)} e de \VUgras{mortero} \textsuperscript{(65)}.

%%K t05-01-20 p
\textsc{Le oncle}. --- Justo, vide ille homine, con un grande cappello, sur su
\VUgras{carro} \textsuperscript{(136)}. Ille es un \VUgras{carrettero}
\textsuperscript{(135)}. Cata die ille apporta hic \VUgras{blocos de petra}
\textsuperscript{(60)}. Ille discarga su vehiculo in le corte, e le
\VUgras{talia-petras} \textsuperscript{(83)} talia le blocos e los seca con lor
\VUgras{serra a petra} \textsuperscript{(85)}. Un \VUgras{manobrero}
\textsuperscript{(146)} los porta al muratero qui los pone in loco.

%%K t05-01-21 p
Intertanto, le \VUgras{mesclator} \textsuperscript{(87)} prepara le mortero.
Ille ha besonio de \VUgras{sablo} \textsuperscript{(62)}, de \VUgras{calce}
\textsuperscript{(63)} e de \VUgras{aqua} \textsuperscript{(64)}. Le calce es
presso ille in un \VUgras{sacco} \textsuperscript{(71)}, un \VUgras{servitor de
muratero} \textsuperscript{(111)} le apporta le sablo sur un \VUgras{carretta a
mano} \textsuperscript{(112)} e le \VUgras{apprentisse}
\textsuperscript{(113)} es al pumpa (58). Ille plena un \VUgras{situla}
\textsuperscript{(114)}. Le mortero es tosto preste. Le \VUgras{portator de
mortero} \textsuperscript{(107)} lo prende e lo porta in alto, per le scala, sur
le scafalage, ubi labora un muratero qui fini iste latere del casa.

%%K t05-01-22 p
E ora, como es nominate le obrero qui labora le \VUgras{tecto}
\textsuperscript{(31)}?

%%K t05-01-23 p
\textsc{Johannes}. --- Ille es un \VUgras{copritor de tectos}
\textsuperscript{(118)}.

%%K t05-tit-4 sub
\VUsaut{3.0mm}
\VUpk{10.2pt}{40}{Le tecto del casa.}
\VUsaut{3.0mm}

%%K t05-01-24 p
S\textsuperscript{o} \textsc{Blanco}. --- Si, ma primo veni le \VUgras{carpentero}
\textsuperscript{(115)}.

%%K t05-01-25 p
Le carpentero labora le \VUgras{carpenteria} \textsuperscript{(35)}. Ille
assembla le \VUgras{traves} \textsuperscript{(37)} per \VUgras{chevilias}
\textsuperscript{(117)}. Ille obreros, illac in alto, con lor large pantalones
de veluto, es carpenteros. Le un tene un parve trave, e le altere seca un
chevilia con su \VUgras{hachetta} \textsuperscript{(116)}. Ille qui es presso le
\VUgras{camino} \textsuperscript{(41)} es le copritor de tectos. Ille inchoda
lattas sur le (*) trabes, e \VUgras{ardesias} \textsuperscript{(66)} sur le
lattas. Alcun vices ille pone tegulas.

%%K t05-noto-1 noto
\VUnotes{143.2mm}{%
(*) \VUgras{Balk-o} = G. \textit{Balken}, A. \textit{joist}, F. \textit{solive},
I. \textit{travicello}, R. \textit{balka}, E. Port. \textit{viga}. Radice
technic usque ora non adoptate, ma proponibile pro traducer le senso del parola
F. \textit{solive}, que non es un trabe F. \textit{poutre}, ni un parve trave.
Illo es un barra de ligno quadrate que supporta un solo e un plafon.%
}

%%K t05-01-26 p
Le tecto del stabulo es \VUgras{coperite de tegulas} \textsuperscript{(55)},
illo del casa es \VUgras{coperite de ardesias} \textsuperscript{(32)} e illo del
veranda es de \VUgras{zinc} \textsuperscript{(24)}.

%%K t05-01-27 p
\textsc{Johannes}. --- Esque le copritor de tectos labora anque le tectos de
zinc?

%%K t05-01-28 p
S\textsuperscript{o} \textsc{Blanco}. --- No, ma le \VUgras{obrero de zinc}
\textsuperscript{(121)} o le \VUgras{plumbero}, ille qui pone un
\VUgras{lucarna} \textsuperscript{(38)} sur le tecto del casa. Ille pone anque
le \VUgras{tubos de pluvia} \textsuperscript{(43)} e le \VUgras{gutteras del
tecto} \textsuperscript{(42)}. Ille solda le zinc e le latta. Ille fixava le
\VUgras{parafulmine} \textsuperscript{(40)} e le \VUgras{gyroscopio}
\textsuperscript{(39)} sur le \VUgras{fronton} \textsuperscript{(36)}.

%%K t05-01-29 p
\textsc{Johannes}. --- Assi le casa es finite?

%%K t05-tit-5 sub
\VUsaut{3.0mm}
\VUpk{10.2pt}{40}{Le utensiles.}
\VUsaut{3.0mm}

%%K t05-01-30 p
S\textsuperscript{o} \textsc{Blanco}. --- Non completemente. Le
\VUgras{gypsero} \textsuperscript{(99)} construe con \VUgras{briccas}
\textsuperscript{(61)} le parietes que lo divide in diverse cameras. Ille coperi
de \VUgras{gypso} \textsuperscript{(70)} le \VUgras{plafones}
\textsuperscript{(26)} e le muros. Ora ille labora le plafon del
\VUgras{veranda} \textsuperscript{(33)}. Ille ha, como le muratero, un
\VUgras{trulla} \textsuperscript{(100)} e un \VUgras{gaveta}
\textsuperscript{(101)}.

%%K t05-01-31 p
Postea le ebenista labora le portas, le fenestras, le \VUgras{parquetes}
\textsuperscript{(16)} e le \VUgras{persianas} \textsuperscript{(19)}.

%%K t05-01-32 p
Ora ille labora in le corte sur su \VUgras{banco} de travalio
\textsuperscript{(90)}. Ille ha un \VUgras{serra} \textsuperscript{(94)} pro
secar le ligno (69), un \VUgras{plana} \textsuperscript{(91)}, un
\VUgras{morsa} \textsuperscript{(92)}, un \VUgras{cisello}
\textsuperscript{(94 bis)}, un \VUgras{compasso} \textsuperscript{(95)}, e un
\VUgras{martello} de ligno \textsuperscript{(95 bis)}.

%%K t05-01-33 p
\textsc{Johannes}. --- Ha! ma il es plus amusante facer ebenisteria que
muratura. Oncle mie, tu comprara pro me un banco de travalio, un serra e un
plana, proque io tosto laborara un cabana pro Medor.

%%K t05-01-34 p
\textsc{Le oncle}. --- Si, parve puero.

%%K t05-01-35 p
\textsc{Johannes}. --- Quanto contente io es! E qui es ille qui sta presso le
porta de entrata?

%%K t05-tit-6 sub
\VUsaut{3.0mm}
\VUpk{10.2pt}{40}{Le pictura.}
\VUsaut{3.0mm}

%%K t05-01-36 p
S\textsuperscript{o} \textsc{Blanco}. --- Ille es un \VUgras{pictor}
\textsuperscript{(102)}. Ille sta sur su scala con un pote de \VUgras{pictura}
\textsuperscript{(103)} e su \VUgras{pincel} \textsuperscript{(104)}. Ille pinge
le ligno del porta de entrata pro un plus longe conservation.

%%K t05-01-37 p
Ille anque colla le papiros in le appartamentos.

%%K t05-01-38 p
Le \VUgras{tapissero} \textsuperscript{(107)} qui es sur un \VUgras{scala}
\textsuperscript{(106)}, sur le \VUgras{balcon} \textsuperscript{(28)}, pone un
\VUgras{store} \textsuperscript{(29)}.

%%K t05-01-39 p
Le obrero stante presso le grande fenestra es le \VUgras{vitrero}
\textsuperscript{(96)}. Ille fixa le \VUgras{vitros} \textsuperscript{(18)} per
\VUgras{mastico} \textsuperscript{(73)}. Ille altere obrero, genuflexe presso le
porta del cellario es un \VUgras{serrallero} \textsuperscript{(97)}. Ille pone
un \VUgras{serratura} \textsuperscript{(6)}. Su \VUgras{lima}
\textsuperscript{(98)} es presso ille.

%%K t05-01-40 p
\textsc{Johannes}. --- Esque es anque un serrallero, ille obrero qui batte con
su \VUgras{martello} \textsuperscript{(84)} sur un ferramento?

%%K t05-01-41 p
S\textsuperscript{o} \textsc{Blanco}. --- Plus exactemente, ille es un ferrero
(125). Ille forgia \VUgras{ferramentos} \textsuperscript{(67)} pro le
\VUgras{electricista} \textsuperscript{(124)} qui es sur le tecto del
\VUgras{remissa} \textsuperscript{(51)}. Ille ha un \VUgras{forgia}
\textsuperscript{(126)}, un \VUgras{incude} \textsuperscript{(127)} e un
\VUgras{tenalia} \textsuperscript{(128)}.

%%K t05-01-42 p
Le electricista fixa le \VUgras{filo} de \VUgras{cupro} \textsuperscript{(44)}
del telephono.

%%K t05-tit-7 sub
\VUpk{10.2pt}{40}{Le jardinage.}
\VUsaut{3.0mm}

%%K t05-01-43 p
\textsc{Le oncle}. --- In le fundo del \VUgras{corte} \textsuperscript{(45)},
presso le \VUgras{grillia} \textsuperscript{(47)} e le \VUgras{portal}
\textsuperscript{(48)} tu vide le \VUgras{jardinero} \textsuperscript{(129)}.
Ille prepara le \VUgras{parve jardin} \textsuperscript{(49)}. Ille fodeva le
terra con su \VUgras{pala} \textsuperscript{(131)}, ille semina semines e rastra
con le \VUgras{rastro} \textsuperscript{(130)}. Postea, con su
\VUgras{regatorio} \textsuperscript{(132)}, ille aquara le \VUgras{plattabandas}
\textsuperscript{(150)} e le plantas crescera.

%%K t05-01-44 p
Le \VUgras{palafrenero} \textsuperscript{(133)} qui justo curava le cavallo e le
dava alimento netta le \VUgras{vehiculo} \textsuperscript{(52)} con un spongia,
un \VUgras{pumpa a mano} \textsuperscript{(134)} e un situla de aqua. Postea
ille lo reintroducera in le remissa e claudera le porta con le spisse
\VUgras{pessulo} \textsuperscript{(53)}.

%%K t05-01-45 p
Ecce te contente, io crede, tu habeva bastante explicationes.

%%K t05-01-46 p
\textsc{Johannes}. --- Si, oncle mie, ma con placer io viderea le interior del
casa.

%%K t05-01-47 p
\textsc{Le oncle}. --- Isto habera loco deman. Senior Rufo te explicara le plano
e te indicara le nomine de cata camera.

%%K t05-tit-8 sub
\VUsaut{3.0mm}
\VUpk{10.2pt}{40}{Disposition interne del casa.}
\VUsaut{2.4mm}

%%K t05-apar-2 apar
{\centering\textit{(Vide le plano.)}\par}
\VUsaut{2.4mm}

%%K t05-01-48 p
\textsc{Le architecto}. --- Veni, Johannes, io te facera tosto explorar le
diverse partes del casa.

%%K t05-01-49 p
Entramus le \VUgras{terreno} \textsuperscript{(7)}. Ecce le button del
\VUgras{campanetta} \textsuperscript{(11)}. Nos ascende le tres \VUgras{gradines}
\textsuperscript{(14)} del \VUgras{perron} \textsuperscript{(9)}, nos transversa
le \VUgras{limine} \textsuperscript{(10)} del \VUgras{porta}
\textsuperscript{(8)} de entrata e ecce que nos es al pede del \VUgras{scala}
\textsuperscript{(13)}.

%%K t05-01-50 p
\textsc{Johannes}. --- Isto es le corridor.

%%K t05-01-51 p
\textsc{Le architecto}. --- No, isto es le \VUgras{vestibulo}
\textsuperscript{(12)}. Le \VUgras{corridor} \textsuperscript{(a)} es al
extremitate. Al dextra, ecce le \VUgras{salon} \textsuperscript{(b)} e le
\VUgras{sala a mangiar} \textsuperscript{(c)}, opposite al sala a mangiar es le
\VUgras{cocina} \textsuperscript{(d)} e le \VUgras{officio}
\textsuperscript{(e)}, postea ante illo le \VUgras{gabinetto de travalio}
\textsuperscript{(f)}. Le \VUgras{veranda} \textsuperscript{(23)} con su
\VUgras{porta vitrate} \textsuperscript{(25)} es presso le salon.

%%K t05-01-52 p
Nos ascendera tosto le scala. Appoia te al \VUgras{rampa}
\textsuperscript{(15)} e conta le gradines. Illos es dece-quatro. Ecce le
\VUgras{planetto} \textsuperscript{(m)}, nos es sur le prime \VUgras{etage}
\textsuperscript{(27)}.

%%K t05-tit-9 sub
\VUsaut{3.0mm}
\VUpk{10.2pt}{40}{Le prime etage.}
\VUsaut{3.0mm}

%%K t05-01-53 p
Opposite al scala, il ha le \VUgras{latrina} \textsuperscript{(20)} e le
\VUgras{gabinetto de toilette} \textsuperscript{(j)}. Al dextra un belle
\VUgras{camera de dormir} \textsuperscript{(g)} con duo fenestras al
\VUgras{facade} \textsuperscript{(1)} del casa. Al sinistra, il ha le \VUgras{sala
de banio} \textsuperscript{(1)} e le \VUgras{camera pro amicos}
\textsuperscript{(i)} con un grande \VUgras{alcova} \textsuperscript{(h)}:
Presso le camera de dormir es le \VUgras{camera pro le infantes}
\textsuperscript{(k)}. Illo se trova presso un vaste \VUgras{terrassa}
\textsuperscript{(21)}. Sub iste terrassa il ha un altere latrina (20) e, al
muro, ecce le \VUgras{campana} \textsuperscript{(22)} que sona pro le cena.

%%K t05-01-54 p
\textsc{Johannes}. --- Vamos sur le \VUgras{secunde etage}.

%%K t05-01-55 p
\textsc{Le architecto}. --- Il non existe un secunde etage. Sub le tecto il ha
solmente \VUgras{mansardas} \textsuperscript{(33)} e le granario (34). Nos
finiva le exploration, nos pote redescender.

%%K t05-01-56 p
\textsc{Johannes}. --- Gratias, senior, isto es multo interessante. Io
racontara tosto a mi patre toto lo que io videva.
\VUsaut{4.0mm}
\VUfilet{22mm}
